\documentclass[12pt]{article}

\usepackage[spanish]{babel}
\usepackage[utf8]{inputenc}
\usepackage[T1]{fontenc}
\usepackage[margin=2.5cm]{geometry}
\usepackage{setspace}

\onehalfspacing

\begin{document}

\section{Introducción}
La ciberseguridad en el ámbito militar surge como una respuesta necesaria ante nuevas amenazas que no requieren ejercicios tradicionales ni armamento convencional, sino conocimientos técnicos y acceso a la red. Ataques cibernéticos como los sufridos por Estonia, Georgia o incluso Estados Unidos demuestran que un conflicto puede comenzar sin disparar una sola bala. Por ello, la ciberseguridad se ha convertido en un elemento clave para la defensa nacional e internacional, obligando a Estados y organizaciones como la OTAN a replantear sus doctrinas de seguridad, sus estructuras de defensa y sus mecanismos de cooperación frente a los ciberataques.

\section{Justificación}
El estudio de la ciberseguridad en el ámbito militar es fundamental debido al creciente impacto que los ciberataques pueden tener sobre la estabilidad, la seguridad y la soberanía de los Estados. La alta dependencia tecnológica de las Fuerzas Armadas hace que un ataque a sus sistemas de información pueda paralizar operaciones, comprometer información clasificada o incluso afectar infraestructuras críticas del país. Esto convierte a la ciberseguridad en una prioridad estratégica, al mismo nivel que otras amenazas tradicionales como el terrorismo o la proliferación de armas de destrucción masiva. Analizar este tema permite comprender la necesidad de desarrollar capacidades de ciber defensa.

\section{Marco Teórico}
La ciberseguridad en el ámbito militar se refiere al conjunto de medidas destinadas a proteger los sistemas de información y redes utilizadas por las Fuerzas Armadas. Estos sistemas son esenciales para el mando, control, comunicaciones e inteligencia, por lo que su protección es clave para garantizar la seguridad nacional. El ciberespacio se ha convertido en un nuevo dominio de conflicto, donde las amenazas digitales pueden afectar directamente la capacidad operativa de los Estados.

\section{Antecedentes}
El uso de tecnologías de la información en el ámbito militar aumentó significativamente a finales del siglo XX, incrementando también los riesgos de ataques cibernéticos. Un antecedente clave fue el ataque a Estonia en 2007, que evidenció la vulnerabilidad de las infraestructuras digitales. Estos hechos impulsaron a los Estados y a organizaciones como la OTAN a desarrollar estrategias y capacidades de ciber defensa.

\section{Base Teórica}
La base teórica de la ciberseguridad militar se fundamenta en la protección de la información bajo los principios de confidencialidad, integridad y disponibilidad. Además, retoma conceptos estratégicos clásicos que destacan la importancia de debilitar al enemigo sin un enfrentamiento directo, aplicados hoy al ciberespacio.

\section{Objetivo General}
El objetivo general busca comprender de manera integral cómo y por qué la ciberseguridad se ha vuelto un pilar estratégico en el ámbito militar. No se trata solo de proteger computadoras, sino de asegurar la capacidad de defensa de un Estado frente a amenazas que ya no necesitan armas físicas.

Este objetivo apunta a:
\begin{itemize}
    \item Reconocer al ciberespacio como un nuevo campo de batalla.
    \item Entender que los ciberataques pueden afectar decisiones militares, infraestructuras críticas y estabilidad nacional.
    \item Analizar las respuestas estratégicas de los Estados y de organismos internacionales como la OTAN ante estos riesgos.
\end{itemize}

\section{Objetivos Específicos}

\begin{enumerate}
    \item Identificar las amenazas y riesgos en el ciberespacio.  
    Este objetivo pretende reconocer qué tipo de ataques existen, como ciberataques, terrorismo digital o guerras cibernéticas, y cómo estos pueden afectar directamente a las Fuerzas Armadas, desde el robo de información hasta la paralización de sistemas militares.

    \item Analizar el ciberespacio como nuevo dominio de conflicto.  
    Se busca explicar por qué el ciberespacio ya se considera un dominio militar más, al mismo nivel que la tierra, el mar y el aire, comprendiendo que los conflictos modernos no siempre son visibles, pero sí devastadores.

    \item Describir las operaciones cibernéticas militares.  
    Este objetivo se enfoca en explicar cómo funcionan las operaciones de:
    \begin{itemize}
        \item Ciber defensa (protección de sistemas).
        \item Ciber explotación (obtención de información).
        \item Ciberataque (daño o sabotaje al enemigo).
    \end{itemize}

    \item Examinar la evolución de las estrategias de la OTAN.  
    Busca analizar cómo la OTAN ha tenido que adaptar su concepto estratégico ante las nuevas amenazas digitales, incorporando la ciberseguridad como un tema prioritario.

    \item Analizar casos reales de ciberataques.  
    Este objetivo pretende usar ejemplos reales, como los casos de Estonia o Georgia, para demostrar que los ciberataques no son teoría, sino hechos reales con consecuencias políticas, económicas y militares.

    \item Resaltar la importancia de la cooperación y la formación.  
    Se busca destacar que la ciberseguridad no depende solo de la tecnología, sino también de personal capacitado, cooperación internacional y concientización y adiestramiento constante.
\end{enumerate}

\section{Conclusión}
La ciberseguridad en el ámbito militar se ha consolidado como un elemento esencial para la defensa y seguridad de los Estados en la actualidad. El avance tecnológico y la creciente dependencia de los sistemas digitales han transformado la manera en que se planifican y ejecutan las operaciones militares, haciendo del ciberespacio un nuevo dominio estratégico de conflicto. En conclusión, la ciberseguridad militar no solo implica el uso de tecnologías avanzadas, sino también la capacitación internacional y el desarrollo de marcos normativos adecuados, frente a un entorno digital cada vez más complejo y vulnerable.

\end{document}